\phantomsection
\color{unidarkgreen}\section*{\centering{\large{ПЕРЕЧЕНЬ СОКРАЩЕНИЙ}}}
\addcontentsline{toc}{section}{Перечень сокращений}

% ===== НАСТРОЙКИ =====
\newlength{\abbrgap}
\setlength{\abbrgap}{1.5em}   % <-- расстояние между "--" и расшифровкой (меняйте здесь)

% Автоматическая ширина метки: самое длинное сокращение + тире + пробел + запас
\newlength{\abbrwidth}
\settowidth{\abbrwidth}{АСУ~ТП}   % самое длинное сокращение
\addtolength{\abbrwidth}{3.5em}% запас под "-- " и отступ
% ======================

{\color{white}\fontsize{0.1pt}{0.1pt}\selectfont<begABBRS>}
\color{black}

\begin{list}{}%
{%
  \setlength{\labelwidth}{\abbrwidth}%
  \setlength{\labelsep}{0pt}%      <-- теперь весь отступ задаётся через \abbrgap в метке
  \setlength{\leftmargin}{\dimexpr\labelwidth+\labelsep\relax}%
  \setlength{\itemindent}{0pt}%
  \setlength{\parsep}{0pt}%
  \setlength{\itemsep}{0pt}
  \renewcommand{\makelabel}[1]{#1 --\hspace{\abbrgap}}%   <-- метка = сокращение + "--" + пробел
}
\item[GOOSE] Generic Object-Oriented Substation Event;
\item[IRF] Internal Relay Fault;
\item[PPS] Pulse Per Second;
\item[PTP] Precision Time Protocol (протокол точного времени);
\item[SNTP] Simple Network Time Protocol (простой сетевой протокол времени);
\item[АПВ] автоматическое повторное включение;
\item[АСУ~ТП] автоматизированная система управления технологическими процессами;
\item[АУВ] автоматика управления выключателем;
\item[АЦП] аналого-цифровой преобразователь;
\item[БИ] блок испытательный;
\item[БКУ] блок команд управления;
\item[БНН] блокировка при неисправности цепей напряжения;
\item[БНТ] блокировка при бросках намагничивающего тока;
\item[ВВ] высоковольтный ввод;
\item[ВН] высшее напряжение;
\item[ВЧ] высокочастотный;
\item[ВЧПП] высокочастотный приемо-передатчик;
\item[ДЗ] дистанционная защита;
\item[ДЗШ] дифференциальная защита шин;
\item[ДО] дочерние общества;
\item[ДТЗ] дифференциальная токовая защита;
\item[ДТЗт] дифференциальная токовая защита с торможением;
\item[ДТО] дифференциальная токовая отсечка;
\item[ДФЗ] дифференциально-фазная защита;
\item[ЕЭС] единая энергетическая система;
\item[ЗДЗ] защита от дуговых замыканий;
\item[ЗИП] запасные части, инструменты и принадлежности;
\item[ЗНР] защита от неполнофазного режима;
\item[ЗНФ] защита от непереключения фаз;
\item[ИО] измерительный орган / избирательный орган;
\item[ИЧМ] интерфейс человек-машина;
\item[КЗ] короткое замыкание;
\item[КИ] контроль изоляции;
\item[КО] компенсационная обмотка;
\item[КОН] контроль отсутствия напряжения;
\item[КС] контроль синхронизма / контрольная сумма;
\item[КСВ] контроль силового выключателя;
\item[КСН] контроль синхронизма и напряжения;
\item[КЦТ] контроль цепей тока;
\item[ЛО] логика отключения;
\item[ЛЭП] линия электропередачи;
\item[МТЗ] максимальная токовая защита;
\item[МЭК] международная электротехническая комиссия;
\item[НВЧЗ] направленная высокочастотная защита;
\item[НН] низшее напряжение;
\item[ННЭ] наличие напряжения на энергообъекте;
\item[ОВ] оперативный вывод / обходной выключатель;
\item[ОЗ] основная защита;
\item[ОЗУ] оперативное запоминающее устройство;
\item[ОНЭ] отсутствие напряжения на энергообъекте;
\item[ОТ] оперативный ток;
\item[ОУ] оперативное ускорение;
\item[ПА] противоаварийная автоматика;
\item[ПАВ] полуавтоматическое включение;
\item[ПЗУ] постоянное запоминающее устройство;
\item[ПК] предохранительный клапан;
\item[ПО] пусковой орган / программное обеспечение;
\item[ПРД] передатчик;
\item[ПРМ] приемник;
\item[РД] реле давления;
\item[РЗ] релейная защита;
\item[РЗА] релейная защита и автоматика;
\item[РКВ] реле команды «включить»;
\item[РКО] реле команды «отключить»;
\item[РС] регистрация событий / реле сопротивления;
\item[РТ] реле тока;
\item[РУ] распределительное устройство;
\item[РФ] реле фиксации;
\item[РФК] реле фиксации команд;
\item[РФП] реле фиксации включенного положения;
\item[РЭ] руководство по эксплуатации;
\item[СО] система охлаждения;
\item[ССПИ] система сбора и передачи информации;
\item[ТЗНП] токовая защита нулевой последовательности;
\item[ТЗОП] токовая защита обратной последовательности;
\item[ТК] токовый контроль / текущий контроль;
\item[ТМПд] трансформатор масляный подмагничивания динамических режимов;
\item[ТМПо] трансформатор масляный подмагничивания основной;
\item[ТМПр] трансформатор масляный подмагничивания резервный;
\item[ТН] трансформатор напряжения;
\item[ТР] технический регламент;
\item[ТС] технологическая сигнализация;
\item[ТТ] трансформатор тока;
\item[ТУ] телеуправление / телеускорение / технические условия;
\item[УВ] управление выключателем / управляющее воздействие;
\item[УРОВ] устройство резервирования при отказе выключателя;
\item[УС] улавливание синхронизма;
\item[УШР] управляемый шунтирующий реактор;
\item[ФБ] функциональный блок;
\item[ФПВ] фиксация положения выключателя;
\item[ФСУ] функциональная схема устройства;
\item[ЦВ] цепи включения;
\item[ЦО] цепи отключения;
\item[ЦОС] цифровая обработка сигналов;
\item[ЦП] центральный процессор;
\item[ЦПС] цифровая подстанция;
\item[ЦУ] цепи управления;
\item[ШП] шинки питания;
\item[ШЭТ] шкаф электротехнический типовой;
\item[ЭМВ] электромагнит включения;
\item[ЭМО] электромагнит отключения;
\item[ЭМУ] электромагнит управления.
{\color{white}\fontsize{0.1pt}{0.1pt}\selectfont<endABBRS>}
\end{list}
